\documentclass[11pt]{article}
\pdfmapfile{=lm.map}\pdfmapfile{+cm.map}\pdfmapfile{+symbols.map}
\usepackage[T1]{fontenc}
\usepackage{lmodern,amsmath,amssymb,amsthm,mathtools,booktabs,array}
\usepackage[margin=0.9in]{geometry}
\usepackage{xcolor,microtype,fancyhdr,enumitem,graphicx}
\usepackage[hidelinks]{hyperref}
\definecolor{ink}{HTML}{18334B}
\pagestyle{fancy}\fancyhf{}\fancyhead[L]{\small MOC: categorical follow-up}
\fancyhead[R]{\small 17 October 2026}\fancyfoot[C]{\thepage}
\setlength{\headheight}{14pt}\setlength{\parindent}{0pt}\setlength{\parskip}{6pt}
\newtheorem{theorem}{Theorem}[section]\newtheorem{proposition}[theorem]{Proposition}
\newtheorem{corollary}[theorem]{Corollary}
\newcommand{\C}{\mathbb C}\newcommand{\R}{\mathbb R}\newcommand{\conv}{\operatorname{conv}}
\newcommand{\diag}{\operatorname{diag}}\newcommand{\Rea}{\operatorname{Re}}
\newcommand{\Ima}{\operatorname{Im}}\newcommand{\rank}{\operatorname{rank}}
\newcommand{\id}{\mathrm{id}}
\begin{document}
\begin{center}
{\LARGE\bfseries\color{ink} MOC: a categorical follow-up}\\[7pt]
{\large Exact composition, coherent cycles, and an unbounded spectral family}\\[7pt]
17 October 2026
\end{center}

\textbf{Outcome.} This round establishes additional restricted MOC theorems and exact structural results. It does not prove or refute unrestricted complex MOC. Six parallel research tracks used nLab as a conceptual index, followed primary mathematical sources, developed proofs, and independently checked the strongest claims. No publication-priority claim is made.

For $a,b\in\C^n$ and $U\in U(n)$, write
\[
d_{a,b}(U)=\det(D_a+UD_bU^*),\qquad
z_\pi(a,b)=\prod_i(a_i+b_{\pi(i)}).
\]
MOC asserts $d_{a,b}(U)\in\conv\{z_\pi(a,b):\pi\in S_n\}$ for every choice of these data.

\begin{center}\begin{tabular}{p{0.27\textwidth}p{0.61\textwidth}}
\toprule Result & Proved scope\\\midrule
An unbounded complex family & The preceding interval $|u|\le1/20$ is extended to \emph{every real $u$}, for all $U(4)$; admissible common projective transforms are also covered.\\
Exact categorical composition & For two fixed coordinate partitions, squared-minor kernels compose for all legal block unitaries \emph{if and only if} their bipartite incidence multigraph is a forest.\\
Positive laws on cycles & Explicit corrected permutation weights give an all-dimensional phase region for two matching layers. Balanced gates are covered at \emph{every phase}.\\
Phase behavior & The independent phase count is the graph cycle rank. A single cycle has an exact ellipse orbit; a fixed $U(6)$ phase family is certified rationally.\\
Compatible stochastic lifts & A composition- and tensor-closed incidence category captures exactly the three-line singleton input theorem.\\\bottomrule
\end{tabular}\end{center}

The proofs below distinguish coherent unitary composition, stochastic composition, and compatibility of coupled marginals. These distinctions supply both positive theorems and exact tests of proposed extensions. Ordinary categorical terminology is not being used as an extra axiom.

\newpage
\section{The categorical objects and the positivity requirement}
Let $H_n=(\C^2)^{\otimes n}$, with Boolean basis $|I\rangle$, $I\subseteq[n]$. Define
\[
R(U)_{I,J}=|\det U[I,J]|^2\quad(|I|=|J|),
\]
and zero between different weights. Every weight block is doubly stochastic by Cauchy--Binet. The exterior compression
\[
J(v_1\otimes\cdots\otimes v_n)
 =(e_1\otimes v_1)\wedge\cdots\wedge(e_n\otimes v_n),
\qquad
R(U)=J^*\frac{\bigwedge^n(U\otimes I_2)}{\det U}J
\]
shows that $R(U)$ commutes with the diagonal $GL_2(\C)$ action.
Schur--Weyl duality places it in the two-row permutation quotient, equivalently the classical Temperley--Lieb algebra at loop parameter two. This exact background is already in N059 of the supplied volume.

For product input $u=\bigotimes_j(1,b_j)^T$ and product covector $w=\bigotimes_i(a_i,1)$,
\[
wR(U)u=d_{a,b}(U).
\]
More generally $w_i(u_j)$ is a rank-at-most-two matrix. Generic row and column rescaling reduces it to the form $a_i+b_j$; continuity handles the remaining data. Thus unrestricted MOC is exactly a \emph{two-sided product-test} convexity assertion.


\subsection{An exact identity used by the certificate checks}
For $U$ unitary, let
\[
q_\pi(U)=\Rea\left(\frac{\operatorname{sgn}(\pi)\prod_iU_{i,\pi(i)}}{\det U}\right).
\]
Laplace expansion and the complementary-minor identity give
\[
\sum_\pi q_\pi=1,\qquad
|\det U[I,J]|^2=\sum_{\pi:\,\pi(I)=J}q_\pi,\qquad
d_{a,b}(U)=\sum_\pi q_\pi z_\pi(a,b).
\tag{1}
\]
The $q_\pi$ need not be positive. Equation (1) means that a linear identity between common-minor functions, checked on every permutation, holds on every unitary. It is the bridge used below to turn polynomial coefficient checks into all-unitary support inequalities.


The exact averaged law
\[
\mathbb E_D R(UDV)=R(U)R(V)
\tag{2}
\]
holds for independent uniform diagonal phases $D$. Expanding each minor and averaging its characters proves it. Equation (2), already in N059, inserts dephasing. It does not turn the coherent product $UV$ into an ordinary Markov composition. Markov categories provide the correct language for stochastic gluing and conditional dependence \cite{fritz}; the following theorem identifies where coherent composition itself agrees with it.

\section{A complete classification of exact block composition}
Let $\mathcal P,\mathcal Q$ be partitions of $[n]$. A unitary respects a partition if it is block diagonal on its blocks. Construct the bipartite incidence \emph{multigraph} $G(\mathcal P,\mathcal Q)$: left vertices are $\mathcal P$-blocks, right vertices are $\mathcal Q$-blocks, and coordinate $i$ is an edge between its two blocks. Two coordinates in the same intersection give parallel edges, which constitute a two-edge cycle.

\begin{theorem}\label{forest}
The following statements are equivalent:
\begin{enumerate}[nosep]
\item $G(\mathcal P,\mathcal Q)$ is a forest.
\item $R(UV)=R(U)R(V)$ for every $U$ respecting $\mathcal P$ and $V$ respecting $\mathcal Q$.
\end{enumerate}
If the individual factors have common positive permutation laws, the composite has their product law.
\end{theorem}
\begin{proof}
For fixed equal-size $I,J$, Cauchy--Binet gives
\[
\det(UV)[I,J]=\sum_K \det U[I,K]\det V[K,J].
\]
A nonzero term must obey
\[
|K\cap P|=|I\cap P|\quad(P\in\mathcal P),\qquad
|K\cap Q|=|J\cap Q|\quad(Q\in\mathcal Q).
\tag{3}
\]
Regard $K$ as a zero-one edge set with specified degree at every graph vertex. In a forest it is unique if it exists: subtract two solutions and eliminate leaves; the difference on each leaf edge is zero. There is at most one nonzero amplitude, so squaring the Cauchy--Binet sum gives precisely $R(UV)=R(U)R(V)$, simultaneously at every weight.

Conversely, choose any simple even cycle, allowing a parallel-edge cycle. At each left and right cycle vertex, mix its two incident coordinate labels with a real two-site rotation having all entries nonzero; act as the identity on all other labels. Let $K_0,K_1$ be the two alternating cycle edge sets and take $I=J=K_0$. The only nonzero intermediate sets are $K_0,K_1$. Their amplitudes are nonzero real numbers $a,b$. The composition defect is
\[
(a+b)^2-a^2-b^2=2ab\ne0.
\]
The rational rotation with cosine $3/5$ and sine $4/5$ suffices. Thus a cycle excludes \emph{uniform} exact composition.

Finally, if $R(U)=\sum_p\lambda_pR(P_p)$ and $R(V)=\sum_q\mu_qR(P_q)$, the forest identity yields
$R(UV)=\sum_{p,q}\lambda_p\mu_qR(P_pP_q)$. These weights are nonnegative and sum to one.
\end{proof}

\begin{corollary}\label{hypertree}
Build a circuit of unitary gates such that each new support meets every previously connected component in at most one coordinate. If every gate has a common permutation law, so does the entire circuit. In particular arbitrary complex two-site and three-site gates satisfying this condition give MOC for \emph{all complex spectra in every dimension}.
\end{corollary}
\begin{proof}
At each step the previous unitary is block diagonal on its connected components. Its partition and the new support partition have forest incidence, so Theorem~\ref{forest} applies inductively. Every $U(3)$ has a common law: Birkhoff decomposition handles its squared entries, complementary minors give the same law at weight two, and weights zero and three are one.
\end{proof}

The support condition is equivalently acyclicity of the bipartite graph of sites and gate supports. Pairwise intersections of size at most one alone are insufficient: three supports can still form a cycle. This theorem permits arbitrary internal three-site interference before blocks are joined, extending the unrepeated two-site forest construction in the supplied volume.

For example, embed
\[
H_1=\frac19\begin{pmatrix}7&-4&-4\\-4&1&-8\\-4&-8&1\end{pmatrix},
\qquad
H_2=\frac17\begin{pmatrix}6&-2&-3\\-2&3&-6\\-3&-6&-2\end{pmatrix}
\]
on $\{1,2,3\}$ and $\{3,4,5\}$ respectively. Their product is connected, not a direct sum. The exact verifier constructs 20 positive permutation weights and checks all 252 equal-size squared minors against the same law. The full matrix and weights are printed in the accompanying notes.

\section{Cycles can also have positive laws}
Let $n=2r$, and let $\mathcal P,\mathcal Q$ now be perfect matchings whose union is one alternating cycle. Let $U,V$ be arbitrary complex two-site layers on these matchings. There are $2r$ gates in total. Choose orders $(p_0,p_1)$ and $(q_0,q_1)$ on their pairs.

For a transversal $I$ selecting one site per $\mathcal P$-pair define
$f_{\mathcal P}(I)=(-1)^{\#\{\text{selected second sites}\}}$, and set it to zero on other subsets. Define $f_{\mathcal Q}$ similarly. Let $K_0,K_1$ be the two simultaneous transversals, $I_0$ the first-site $\mathcal P$ transversal, and $J_0$ the first-site $\mathcal Q$ transversal. Put
\[
\kappa=\det U[I_0,K_0]\,\overline{\det U[I_0,K_1]}
        \det V[K_0,J_0]\,\overline{\det V[K_1,J_0]},
\]
\[
\theta=\sum_{j\in K_0}\theta_j-\sum_{j\in K_1}\theta_j,\quad
c=\Rea(e^{i\theta}\kappa),\quad
\sigma=f_{\mathcal P}(K_0)f_{\mathcal Q}(K_0).
\]

\begin{theorem}\label{cycle}
For $D=\diag(e^{i\theta_j})$,
\[
R(UDV)=R(U)R(V)+2c\,f_{\mathcal P}f_{\mathcal Q}^{\,T}.
\tag{4}
\]
Index the classical identity/swap gate choices by $s\in\{0,1\}^{2r}$, and let $\pi_s$ be their resulting permutation in circuit order. If $p_g$ is gate $g$'s swapping probability, set
\[
w_s=\prod_g p_g^{s_g}(1-p_g)^{1-s_g},\qquad
\chi_s=(-1)^{\sum_gs_g}.
\]
Then the exact coherent permutation expansion is
\[
\boxed{\quad R(UDV)=\sum_s\bigl(w_s+\sigma c\chi_s\bigr)R(P_{\pi_s}).\quad}
\tag{5}
\]
It is a common positive law whenever all the displayed weights are nonnegative.
\end{theorem}
\begin{proof}
Two distinct intermediate occupancy solutions on a cycle differ by the alternating cycle flow. This is possible only when every vertex has occupancy one. Hence interference occurs only at weight $r$, for transversal $I,J$, with the two intermediate sets $K_0,K_1$.

In $\det U[I,K_0]\overline{\det U[I,K_1]}$, row regrouping signs cancel. Switching the chosen row in one two-site block changes its factor to its negative by column orthogonality. Thus the product is a fixed scalar times $f_{\mathcal P}(I)$. Row orthogonality gives the corresponding $f_{\mathcal Q}(J)$ factor for $V$. Evaluating at $I_0,J_0$ gives $\kappa$, proving (4).

The signed gate-choice operator for the first matching is
\[
\prod_{e\in\mathcal P}(I-T_e)=f_{\mathcal P}f_{\mathcal P}^{\,T},
\]
because each two-site antisymmetrizer annihilates zero or two occupied sites and has matrix $\left(\begin{smallmatrix}1&-1\\-1&1\end{smallmatrix}\right)$ on one occupied site. The analogous identity holds for $\mathcal Q$. Their inner product has only the two simultaneous transversals as terms, each equal to $\sigma$. Therefore
\[
\sum_s\chi_sR(P_{\pi_s})=2\sigma f_{\mathcal P}f_{\mathcal Q}^{\,T}.
\]
Combine this with (4) and the ordinary Bernoulli law for $R(U)R(V)$ to get (5). Its weights sum to one, since $\sum_s\chi_s=0$.
\end{proof}

Local two-site unitarity also gives
$|\kappa|=\prod_g\sqrt{p_g(1-p_g)}$. A useful sufficient positivity condition is
\[
|c|\le\prod_g\min(p_g,1-p_g).
\tag{6}
\]
For nondeterministic gates, writing $\phi=\theta+\arg\kappa$, this becomes
\[
|\cos\phi|\le
\rho:=\prod_g\sqrt{\frac{\min(p_g,1-p_g)}{\max(p_g,1-p_g)}}.
\tag{7}
\]
There is a positive-width phase region around quadrature for every such choice of biases. If all gates are balanced, $p_g=1/2$, then $\rho=1$ and \emph{every phase} has a common positive law. Disjoint cycle components can be combined by product laws. Deterministic gates have $\kappa=0$ and no defect.


\subsection{A determinant-specific cancellation}
Insert (4) in the full minor expansion of the determinant. Its defect from the positive classical mixture factors exactly as
\[
2c\prod_{(p_0,p_1)\in\mathcal P}(a_{p_1}-a_{p_0})
   \prod_{(q_0,q_1)\in\mathcal Q}(b_{q_0}-b_{q_1}).
\tag{8}
\]
Indeed the sums $\sum_I a_{I^c}f_{\mathcal P}(I)$ and
$\sum_J b_Jf_{\mathcal Q}(J)$ factor pair by pair. Therefore one equal-$a$ pair aligned with $\mathcal P$, or one equal-$b$ pair aligned with $\mathcal Q$, removes the defect at \emph{every phase and every gate bias}. All other spectral values remain arbitrary.

A nonzero-interference rational example uses four gates
$\left(\begin{smallmatrix}3/5&4/5\\-4/5&3/5\end{smallmatrix}\right)$ on
$\mathcal P=\{12,34\}$ and $\mathcal Q=\{13,24\}$, with
$D=\diag((7+24i)/25,1,1,1)$. Here $\rho=81/256>7/25$.
The exact checker verifies all 70 squared-minor entries, 16 nonzero defects, and 16 positive corrected weights, whose minimum is $18873/9765625$. Thus the theorem covers an actual closed circuit with nonzero coherent interference, uniformly over all complex spectra.

\subsection{A further complete family: every $U(2)\otimes U(2)$}
\begin{theorem}\label{tensor}
For arbitrary complex $A,B\in U(2)$, the four-dimensional unitary $A\otimes B$ has a common positive permutation law. Consequently MOC holds for this mixing matrix and every pair of complex four-dimensional spectra.
\end{theorem}
\begin{proof}
Diagonal phases and row/column permutations preserve existence of a common law. Every two-site unitary is diagonally equivalent to
\[
R(p)=\begin{pmatrix}\sqrt{1-p}&\sqrt p\\-\sqrt p&\sqrt{1-p}\end{pmatrix}.
\]
Swapping rows replaces $p$ by $1-p$ up to diagonal signs; exchanging the tensor factors is a permutation of four coordinates. We may therefore assume
$A\otimes B=R(p)\otimes R(q)$ with $0\le p\le q\le1/2$.
Put $x_0=1-2q$, $x_1=2(q-p)$ and $x_2=2p$. These are nonnegative and sum to one. The following six probability laws are explicit:
\begin{center}\small\begin{tabular}{cll}
\toprule Law & Permutation labels & Weights\\\midrule
$\nu_{00}$ & $1234$ & $1$\\
$\nu_{01}$ & $1234;\ 1243,2134$ & $1/2;\ 1/4,1/4$\\
$\nu_{02}$ & $1243,1432,2134,3214$ & $1/4$ each\\
$\nu_{11}$ & $1234,1243,2134,2143$ & $1/4$ each\\
$\nu_{12}$ & $1243,2134;\ 1432,2341,3214,4123$ & $1/4$ each; $1/8$ each\\
$\nu_{22}$ & $1243,1432,2134,2341,3214,3421,4123,4312$ & $1/8$ each\\\bottomrule
\end{tabular}\end{center}
The common law is the positive Bernstein interpolation
\[
\lambda(p,q)=\sum_{0\le i\le j\le2}(2-\delta_{ij})x_ix_j\nu_{ij}.
\tag{9}
\]
Its weights sum to $(x_0+x_1+x_2)^2=1$.

To verify the identity, form $R(p)\otimes R(q)$ using formal variables $c,s,d,e$ with $c^2=1-p$, $s^2=p$, $d^2=1-q$, $e^2=q$. Every squared minor reduces to a polynomial of total degree at most two in $p,q$. The supplied verifier constructs the one- and two-minors symbolically, verifies that every surviving exponent is even, and checks all coefficients against (9). Complementary minors give weights three and four. All 70 equal-cardinality row/column pairs agree exactly over $\mathbb Q[p,q]$. The six displayed control laws therefore prove the identity on the whole triangle, including its boundary; the gauges cover every original $A,B$.
\end{proof}

This is the physical Kronecker product of two $2\times2$ unitaries. It differs from Boolean tensor closure, which represents a direct sum of physical matrices. No assertion about arbitrary Kronecker powers follows from this theorem.

\subsection{Joining positive cyclic pieces}
Any complete common-law family can serve as a seed in Corollary~\ref{hypertree}. Thus positive matching-cycle circuits, the four-site tensor family, and arbitrary three-site unitaries can be joined along single-coordinate incidence forests. The resulting connected all-dimensional families can contain arbitrarily many closed coherent cycles.

For balanced matching cycles, (5) has the transparent form
\[
\lambda_s=2^{-2r}\bigl(1+\sigma\cos\phi\,(-1)^{|s|}\bigr).
\]
The coherent cycle phase becomes a bias between even and odd swap outcomes in a genuine classical probability law.


\section{Cycle phases as graph cohomology}
For general block partitions, add $\alpha_P+\beta_Q$ to the diagonal phase on coordinate edge $e=(P,Q)$. The added block-constant diagonals commute with $U$ and $V$, so $UDV$ changes only by outer diagonal phases. Every squared minor is unchanged. After reversing the sign of the right vertex phases, this is the graph coboundary action.

A spanning forest lets us set all its edge phases to zero. The number of independent remaining phases is
\[
\beta_1(G)=|E|-|V|+\#\{\text{connected components}\}.
\tag{10}
\]
Their invariants are alternating cycle phase sums. This first-cohomology viewpoint is operational: it reduces the parameters without discarding their interference.

When $\beta_1=1$, only one coordinate phase $z=e^{i\theta}$ remains. Each minor of $UD(z)V$ is affine in $z$, by Cauchy--Binet, or because the matrix is a rank-one affine pencil in $z$. Hence its squared modulus has only frequencies $-1,0,1$. For fixed spectra,
\[
d(\theta)=d_0+p\cos\theta+q\sin\theta,\qquad d_0,p,q\in\C.
\tag{11}
\]
This is an ellipse, including degenerate segments and points. The center $d_0$ is the phase average. With several cycles one obtains a finite Laurent polynomial on a $\beta_1$-torus, with each coordinate frequency in $\{-1,0,1\}$.

\begin{proposition}\label{ellipse}
Let $v_j$ be the counterclockwise vertices of a two-dimensional permutation hull, $e_j=v_{j+1}-v_j$, and $[e,w]=\Ima(\overline e\,w)$. The entire ellipse (11) lies in the hull exactly when
\[
h_j:=[e_j,d_0-v_j]\ge
\sqrt{[e_j,p]^2+[e_j,q]^2}\qquad\text{for every edge }j.
\tag{12}
\]
For rational complex data, it suffices to check
$h_j\ge0$ and $h_j^2\ge[e_j,p]^2+[e_j,q]^2$ in rational arithmetic.
\end{proposition}
\begin{proof}
The edge gap along the ellipse is $h_j+[e_j,p]\cos\theta+[e_j,q]\sin\theta$. Its exact minimum is the left side of (12) minus its right side. Intersect the edge half-planes.
\end{proof}

For a concrete $U(6)$ family, use partitions
$\mathcal P=\{125,34,6\}$ and $\mathcal Q=\{146,23,5\}$, whose incidence graph has one cycle. Let $H=H(1,2,2)$ and $K=H(1,2,3)$ be Householder reflections. On the three-site blocks use
$H\diag(1,i,1)K$ and $K\diag((3+4i)/5,1,1)H$; on the two-site blocks use rotations $(3/5,4/5)$ and $(5/13,12/13)$. Call the resulting layers $A,B$, and set
\[
W(\theta)=A\diag(e^{i\theta},1,1,1,1,1)B,
\]
\[
a=(1+i,2-i,-1+2i,3+i,-2-i,2+3i),\quad
b=(-2+i,1-2i,3+3i,-1+i,-2-2i,2-i).
\]
The exact verifier proves all twelve edge inequalities strictly for the entire phase circle. It reconstructs all 720 permutation products and the nondegenerate ellipse using Gaussian rational arithmetic. The smallest center-margin-to-radius ratio is approximately $204.73$; this decimal describes the result, while the proof uses exact squared inequalities.

\section{Full MOC on an unbounded complex spectral family}
\begin{theorem}\label{allreal}
For every real $u$, put
\[
a(u)=(u+i,-u+i,1,2+i),\qquad b=(-1+i,-2-2i,2-i,1+2i).
\]
Then $d_{a(u),b}(U)\in\conv\{z_\pi(a(u),b)\}$ for every $U\in U(4)$.
The spectra are noncollinear; $b$ is simple and $a$ is simple except at $u=0,\pm2$.
\end{theorem}

This extends the previous $|u|\le1/20$ result to the whole real line. The proof is a finite exact certificate, with its full coefficient data and verification-only programs included in the bundle.

\begin{proof}
First take $u\ge0$ and set $x=u/(1+u)$. Each original determinant and matching product has degree at most two in $u$. Thus
\[
Z_\pi(x)=(1-x)^2z_\pi\!\left(\frac{x}{1-x}\right),\qquad
D_U(x)=(1-x)^2d_U\!\left(\frac{x}{1-x}\right)
\]
are polynomials of degree at most two. For $0\le x<1$ the common multiplier is positive, so normalized and original convexity are equivalent.

The saved certificate divides $(0,1)$ into 89 open algebraic cells. Each endpoint is rational or is specified by a rational squarefree polynomial and an exact rational isolating interval. Every cell has a rational interior point. The verifier checks the first endpoint is zero, the last is one, and adjacent endpoints are the same algebraic number, by polynomial gcd and exact root counting. Consequently there are no uncovered gaps.

On each cell, a cyclic list of permutation labels specifies the polygon vertices. For every successive pair $l,r$, the certificate supplies
\[
\Ima\!\left(\overline{Z_r(x)-Z_l(x)}(D_U(x)-Z_l(x))\right)
 =\sum_t\lambda_{lr,t}(x)|\det U[I_t,J_t]|^2,
\tag{13}
\]
where minors have size one or two and every $\lambda$ is a rational polynomial of degree at most four. The checker independently reconstructs the 24 spectral products and verifies (13) coefficient by coefficient on every permutation. Equation (1) extends the identity to every unitary.

Each $\lambda$ is certified nonnegative throughout its cell. Exact polynomial gcd arithmetic extracts the odd-multiplicity root factor. Sturm counts, evaluated at the exact algebraic endpoints, prove it has no root in the open cell. A positive rational interior value then determines the sign; even-multiplicity zeros cannot change it. Identically zero polynomials are accepted separately. This is interval positivity, not sampling.

The verifier also proves strict positivity of all consecutive turns and all pairwise squared distances between the listed vertices, by excluding \emph{all} roots of those polynomials inside the cell. Every directed edge supports all 24 permutation points by its nonnegative certificate. Thus these are precisely the edges of the convex hull. Equation (13) puts $D_U(x)$ inside all its half-planes.

The check comprises 17 polygon types, 726 facet identities, 8,703 nonnegative coefficient polynomials, and 17,424 permutation identities. All arithmetic in the verification is rational or exact algebraic arithmetic.

At each finite boundary $x_0<1$, use continuity of determinants and vertices and compactness of the finite convex-weight simplex. This includes $x=0$ and eigenvalue collisions. The endpoint $x=1$ represents infinity and is unnecessary for finite $u$. Finally, negative $u$ swaps the first two $a$-labels, preserving both the all-unitary assertion and the permutation hull.
\end{proof}

The verification-only program and its algebraic sign routines were independently read and executed. In particular, the review checked open-interval endpoint conventions, equality of neighboring algebraic endpoints, and the positivity of the normalization.

\subsection{Projective transport of the proved family}
Let $\Delta=\alpha\delta-\beta\gamma\ne0$ and suppose the following denominators are nonzero:
\[
a'_i=\frac{\alpha a_i+\beta}{\gamma a_i+\delta},\qquad
b'_j=\frac{\alpha b_j-\beta}{-\gamma b_j+\delta}.
\]
Then
\[
a'_i+b'_j=
\frac{\Delta(a_i+b_j)}
 {(\gamma a_i+\delta)(-\gamma b_j+\delta)}.
\]
The determinant of $D_aU+UD_b$, divided by $\det U$, and every matching product are multiplied by the same nonzero scalar
\[
\frac{\Delta^n}
 {\prod_i(\gamma a_i+\delta)\prod_j(-\gamma b_j+\delta)}.
\]
Complex scalar multiplication preserves convexity. Thus Theorem~\ref{allreal} includes its entire admissible common projective orbit. It still does not cover arbitrary spectral pairs.

\section{A positive category with the correct incidence compatibility}
Let $\mathcal B_n$ denote normalized nonnegative, weight-preserving Boolean kernels commuting with diagonal $GL_2$. Define $\mathcal H_n\subseteq\mathcal B_n$ by requiring that for every adjacent source inclusion $B\subset C$, the output column laws at $B,C$ admit a coupling supported on $X\subset Y$. This is a finite set of weighted Hall inequalities.

Equivalently, a morphism can carry a stochastic \emph{incidence lift} on adjacent pairs, whose two marginals are its Boolean kernels. Retaining the lift as data makes its composition ordinary stochastic composition.

\begin{theorem}\label{incidence}
$\mathcal H_n$ is convex and closed under composition; the family is closed under tensor products. Every actual $R(U)$ belongs to it. Its defining condition is equivalent to the product-input orbit law on inputs with at most three projective lines, one class occurring once.
\end{theorem}
\begin{proof}
For a composite, first sample an adjacent output pair from the first lift, then sample another from the second lift conditionally on that pair. Nesting is preserved and the marginals multiply. Mixtures mix the lifts. In a tensor product an adjacent input changes in exactly one factor: couple that factor by its incidence lift and sample the unchanged factor identically in both endpoints.

For an actual kernel, write $x$ for the exterior product of the smaller column family and $y=x\wedge v$ for the larger one. If $\mathcal F$ is a family of coordinate subsets, $z=P_{\mathcal F}x$ and $N(\mathcal F)$ its upper neighbors, exterior multiplication by the unit vector $v$ has norm at most one. Hence
\[
\|P_{\mathcal F}x\|^2
 =|\langle z\wedge v,y\rangle|
 \le\|P_{\mathcal F}x\|\,\|P_{N(\mathcal F)}y\|.
\]
These are the weighted Hall cuts; finite max-flow supplies the coupling. This is the finite adjacent-rank stochastic domination argument underlying Lyons' theorem \cite{lyons}.

For the input assertion, a global $GL_2$ map sends two input lines to the coordinate lines. The exceptional factor leaves a vector $\alpha|B\rangle+\beta|C\rangle$ with $B\subset C$ adjacent. One incidence coupling gives
\[
R(\alpha|B\rangle+\beta|C\rangle)
 =\sum_{X\subset Y}\mu_{XY}(\alpha|X\rangle+\beta|Y\rangle).
\]
Every summand is a permutation of the same input. Undo the change of basis. Conversely, the orbit law for $|B\rangle+|C\rangle$ provides an incidence coupling by comparing its two weight blocks.
\end{proof}

This gives a genuine positive gluing structure rather than an assumed positivity theorem for all equivariant maps. It recovers the preceding all-dimensional $(r,s,1)$ spectral theorem and makes its extra axiom explicit. The class does not satisfy unrestricted scalar MOC merely by definition.


\section{Positive flag couplings}
A fixed flag of column subsets $J_0\subset J_1\subset\cdots\subset J_n$ has a joint nested determinantal coupling by successive disintegration. The increments define one random permutation $\sigma$ satisfying
\[
\Pr(\sigma(J_k)=I)=|\det U[I,J_k]|^2.
\]
The same law reproduces \emph{all singleton marginals}: if $J_k\setminus J_{k-1}=\{j\}$, then
\[
\Pr(\sigma(j)=i)
 =\Pr(i\in X_k)-\Pr(i\in X_{k-1})=|U_{ij}|^2.
\]
Complementary flag minors follow from Jacobi. This proves convexity for every linear observable built from those minors and singletons, but does not give a common law across different flags.


\subsection{The unlabelled $(2,2,2)$ problem retains five moments}
Consider interval laws on $X\subset Y$ of sizes two and four in $[6]$, with midpoint law uniform on the two intermediate three-subsets of each interval.

There are 90 interval vertices. Their endpoint-plus-midpoint linear map has rank 43, while the map recording only the three color marginals $(X,Y\setminus X,Y^c)$ has rank 38. The five missing linear moments form
\[
\ker\!\left(L:\R^{\binom63}\longrightarrow\R^{\binom62}\right),\qquad
(Lf)(A)=\sum_{Z\supset A}f(Z).
\]
An explicit basis is supplied by the five noncrossing perfect matchings $M$ on six ordered sites:
\[
h_M(Z)=\prod_{(a,b)\in M}\bigl(1_Z(a)-1_Z(b)\bigr).
\]
The exact verifier proves the ranks over $\mathbb Q$, using independent integer relations for upper bounds and nonzero modular minors for lower bounds. These are linear compatibility statements, not existence of a positive coupling. Coordinate-block cases with one rank-one increment per block are proved by independent adjacent couplings; no universal interval-coupling assertion is made.

\subsection{A positive local repair and an indecomposable neighborhood}
The final research round supplied a positive result within this frontier.
Let $B$ be the three-color marginal map on the 90 interval weights and $H$ the five harmonic moment map. An explicit rational $90\times5$ matrix $Q$ satisfies
\[
BQ=0,\qquad HQ=I_5,\qquad
\max_v\sum_j|Q_{vj}|=1.
\]
Thus a color coupling whose weights are at least $m>0$ admits any harmonic correction $\Delta$ with $\|\Delta\|_\infty\le m$: replace its weights by $w+Q\Delta$. Color marginals stay fixed and all weights remain nonnegative.

In contrast, ordinary six-state trades that fix one anchor of each color and permute three other labels preserve all five harmonic moments. Their span has rank 47, the full kernel of the interval map. They cannot repair a midpoint discrepancy. The larger trades in $Q$ are essential for this construction.

\begin{theorem}\label{local222}
Let
\[
U_0=\frac1{8389}
\begin{pmatrix}
863&3348&-1506&984&1266&7320\\
-2664&1223&-2700&-2046&6882&-1716\\
-1986&1044&-199&7884&936&-1506\\
5496&-4866&-2592&1727&2580&366\\
-5010&-5730&1044&120&359&3348\\
-1812&-264&-7278&-258&-3708&-487
\end{pmatrix}.
\]
For every complex unitary $U$ satisfying
\[
\|U-U_0\|_{\mathrm{op}}\le\frac1{23083105288},
\]
MOC holds for arbitrary $a\in\C^6$ and
$b=(b_1,b_1,b_2,b_2,b_0,b_0)$, with arbitrary complex $b_0,b_1,b_2$.
Other placements of the three pairs use the corresponding column permutation of the center.
\end{theorem}
\begin{proof}
The exact certificate verifies $U_0U_0^T=I$ and supplies 90 strictly positive interval weights with minimum
$m=3/70375321$. Their target is
\[
p_X=|\det U[X,12]|^2,\quad
r_Z=|\det U[Z,123]|^2+|\det U[Z,124]|^2,\quad
q_Y=|\det U[Y,1234]|^2.
\]
These comprise 50 coordinates. For every actual unitary, the six intensity relations
\[
\sum_{Z\ni i}r_Z=\sum_{X\ni i}p_X+\sum_{Y\ni i}q_Y
\]
and the mass relation $\sum_Zr_Z=2\sum_Xp_X$ hold. They follow from the projection intensities $P_{ii}$, $2P_{ii}+W_{ii}$ and $(P+W)_{ii}$, where $P$ uses columns $12$ and $W$ columns $34$. The seven relations are independent, and the interval map has rank 43, so every actual target belongs to its column space.

The certificate selects 43 rows and 43 interval columns forming an invertible matrix with inverse infinity operator norm 82. Any change of the full target in that column space is therefore realized by changing those 43 weights, with magnitude at most 82 times the target's infinity-norm change.

For two unitaries, telescoping the $k$ determinant columns and applying Hadamard's inequality gives
\[
\left||\det U[I,J]|^2-|\det V[I,J]|^2\right|
\le2k\|U-V\|_{\mathrm{op}}.
\]
The target coordinates change by at most $12\|U-U_0\|_{\mathrm{op}}$. The radius is exactly $m/(12\cdot82)$, so repaired weights stay nonnegative. All 50 equations hold because both targets lie in the same rank-43 space.

The interval law represents the canonical product input with column classes $12$, $34$, $56$ on three projective lines. A common $GL_2$ transformation transports it to any three distinct projective input lines; scalar pairing proves the spectral conclusion. Coincident values follow by continuity.
\end{proof}

The first two-column projector of $U_0$ has no zero off-diagonal entry, excluding a nontrivial reducing coordinate partition. Thus this is a certified neighborhood beyond the independent-coordinate-block class. The radius is conservative. The universal $(2,2,2)$ statement for arbitrary $U(6)$ remains unproved.


\section{Verification and scope}
The concrete positive results are Theorems~\ref{forest}, \ref{cycle}, \ref{tensor}, \ref{allreal}, \ref{incidence}, and \ref{local222}, together with the exact phase-family certificate. The report uses nLab's entries on Markov categories, Schur--Weyl duality, determinant/exterior algebra, and stochastic maps as navigation; the probability claims are grounded in primary sources and the finite proofs supplied here.

The archive contains the complete source, detailed proof and review notes, exact JSON certificates, and a verification runner. Its final checks use Python's standard library; no optimizer, floating-point root finder, or sampled positivity assumption enters the exact verifiers. Do not run them with Python's \texttt{-O} option, which disables assertions. Agent review is not formal proof-assistant verification or external peer review.

\textbf{Unrestricted MOC remains unresolved by these arguments.} The categorical approach has supplied exact composition and compatibility theorems, positive coherent-cycle laws, and a larger fully proved spectral family. A universal conclusion would still require a proof covering arbitrary unitary interference and arbitrary complex spectra.

\begin{thebibliography}{9}
\bibitem{fritz} T. Fritz, \emph{A synthetic approach to Markov kernels, conditional independence and theorems on sufficient statistics}, Advances in Mathematics 370 (2020), 107239. \url{https://arxiv.org/abs/1908.07021}.
\bibitem{lyons} R. Lyons, \emph{Determinantal probability measures}, Publications Math\'ematiques de l'IH\'ES 98 (2003), 167--212. Theorem 6.2 and Sections 10--11 distinguish nested couplings from stronger union and complete couplings. \url{https://numdam.org/item/10.1007/s10240-003-0016-0.pdf}.
\bibitem{kenyon} R. Kenyon, \emph{Multideterminantal measures}, arXiv:2501.18349v2 (15 July 2025), especially Theorems 6--7 and Section 6.1. \url{https://arxiv.org/html/2501.18349v2}.
\bibitem{nlab} nLab, navigation entries: \emph{Markov category}, \emph{Schur--Weyl duality}, \emph{determinant}, \emph{stochastic map}. \url{https://ncatlab.org/nlab/show/Markov+category}.
\end{thebibliography}
\end{document}
